\section{Local Hankel analysis}\label{sec:hankel}

The purpose of this section is to analyze the local Hankel
contribution $I_H$ around the singular point $s=1$. Recall that
\[
    F(s,y)=\zeta(s)^{2+2y}H(s,y).
\]
The singular part is carried by the power of $\zeta(s)$; the
remaining factor should be made regular at $s=1$ before the local
Hankel integral is evaluated.

\subsection{The local regular factor}\label{sec:local-regular-factor}

We keep the abbreviation
\[
    \alpha=2+2y.
\]
Define
\begin{equation}\label{eq:q-local-def}
    q(s):=(s-1)\zeta(s),\qquad q(1):=1.
\end{equation}
The pole of $\zeta(s)$ at $s=1$ is simple with residue $1$, so
$q$ is holomorphic near $s=1$ after filling in the removable
singularity. After shrinking the neighborhood of $s=1$ if necessary,
we may assume that $q(s)\neq0$ there. We then choose a single
branch of $\log q(s)$ in this neighborhood, normalized by
\[
    \log q(1)=0.
\]
For $s$ in this neighborhood and $y$ near $0$, define
\begin{equation}\label{eq:G-def}
    G(s,y):=q(s)^{2+2y}H(s,y),
    \qquad
    q(s)^{2+2y}:=\exp\bigl((2+2y)\log q(s)\bigr).
\end{equation}
On the corresponding punctured slit neighborhood of $s=1$, the
branch choices from Section~\ref{sec:perron} give
\begin{equation}\label{eq:F-local-factorization}
    F(s,y)=G(s,y)(s-1)^{-2-2y}.
\end{equation}
Here the branch of $(s-1)^{-2-2y}$ is the local Hankel branch
induced by the contour and branch conventions already fixed in
Section~\ref{sec:contour-decomposition}.

\begin{lemma}[Analyticity and boundedness of the local factor]
\label{lem:G-local-analytic}
Fix $0<\eta<1/2$ and $0<\delta<1/2$. There exist
$r\in(0,\delta)$ and $C_{\eta,\delta}>0$ such that $G(s,y)$ is
holomorphic in $(s,y)$ for $|s-1|<r$ and $|y|<\eta$, and
\[
    |G(s,y)|\leq C_{\eta,\delta}
\]
uniformly for $|s-1|\leq r$ and $|y|\leq\eta$. Moreover, on the
punctured slit neighborhood
\[
    0<|s-1|<r,\qquad s\notin\mathcal S,
\]
we have the local factorization
\[
    F(s,y)=G(s,y)(s-1)^{-2-2y}.
\]
\end{lemma}

\begin{proof}
Apply Proposition~\ref{prop:H-analytic} with this fixed $\delta$.
It gives that the function $H(s,y)$ is
holomorphic in $(s,y)$ on the interior of
$\mathcal R_{\delta,\eta}$ and uniformly bounded on
$\mathcal R_{\delta,\eta}$. Choose $r\in(0,\delta)$. Then
$|s-1|\leq r$ implies $\Real(s)\geq1-r>1-\delta$, so
\[
    H(s,y)\ll_{\eta,\delta}1
    \qquad (|s-1|\leq r,\ |y|\leq\eta).
\]

The function $q(s)=(s-1)\zeta(s)$ is holomorphic at $s=1$ after
setting $q(1)=1$. Since $q(1)=1\neq0$, we shrink $r$, preserving
$r\in(0,\delta)$, so that $q$ has no zeros on $|s-1|\leq r$.
Then a holomorphic branch of $\log q(s)$ exists on $|s-1|<r$,
normalized by $\log q(1)=0$. Since this branch is holomorphic on
a neighborhood of the closed disk $|s-1|\leq r$, it is bounded
there. Therefore
\[
    q(s)^{2+2y}
    =\exp\bigl((2+2y)\log q(s)\bigr)
\]
is uniformly bounded for $|s-1|\leq r$ and $|y|\leq\eta$.
Combining this local bound with the bound for $H(s,y)$ gives
$|G(s,y)|\leq C_{\eta,\delta}$.

Finally, on the punctured slit neighborhood,
\[
    \zeta(s)^{2+2y}
    =q(s)^{2+2y}(s-1)^{-2-2y},
\]
with the branch of $(s-1)^{-2-2y}$ inherited from the local
Hankel contour. Combining this identity with
$F(s,y)=\zeta(s)^{2+2y}H(s,y)$ proves the asserted factorization.
\end{proof}

\begin{lemma}[First-order approximation of the local factor]
\label{lem:G-local-values}
Fix $0<\eta<1/2$ and $0<\delta<1/2$, and let
$r\in(0,\delta)$ be chosen as in Lemma~\ref{lem:G-local-analytic}.
Then
\[
    G(1,y)=H(1,y)
\]
uniformly for $|y|\leq\eta$, and
\[
    G(1,0)=1.
\]
Moreover,
\[
    G(s,y)=G(1,y)+O_{\eta,\delta}(|s-1|)
\]
uniformly for $|s-1|\leq r/2$ and $|y|\leq\eta$.
\end{lemma}

\begin{proof}
Since $q(1)=1$ and the branch is normalized by $\log q(1)=0$, we
have
\[
    G(1,y)=q(1)^{2+2y}H(1,y)=H(1,y).
\]
By Proposition~\ref{prop:H-at-zero}, $H(1,0)=1$. Therefore
$G(1,0)=H(1,0)=1$.

For the first-order estimate, use Lemma~\ref{lem:G-local-analytic}.
Since $G$ is holomorphic in $s$ on $|s-1|<r$ and uniformly bounded
for $|s-1|\leq r$, $|y|\leq\eta$, Cauchy's estimate in the
$s$-variable gives
\[
    \frac{\partial G}{\partial s}(s,y)\ll_{\eta,\delta} 1
\]
uniformly for $|s-1|\leq r/2$ and $|y|\leq\eta$. Therefore
\[
    G(s,y)-G(1,y)
    =\int_1^s \frac{\partial G}{\partial s}(w,y)\,dw
    =O_{\eta,\delta}(|s-1|),
\]
uniformly in the stated range.
\end{proof}

\subsection{Evaluation of the local Hankel contribution}
\label{sec:local-hankel-evaluation}

\begin{proposition}[Local Hankel contribution]\label{prop:local-hankel}
With $\eta$ fixed as in Section~\ref{sec:perron}, put
\[
    T=\exp(\sqrt{\log x}),\qquad \rho:=\frac1{\log x}.
\]
For all sufficiently large $x$, uniformly for $|y|\leq\eta$, we have
\[
    I_H
    =
    xH(1,y)\frac{(\log x)^{1+2y}}{\Gamma(2+2y)}
    +O_\eta\left(x(\log x)^{2\eta}\right).
\]
\end{proposition}

\begin{proof}
Put $\alpha=2+2y$. We apply
Lemmas~\ref{lem:G-local-analytic} and~\ref{lem:G-local-values}
with the fixed strip depth $\delta'$ from
Section~\ref{sec:contour-decomposition}; since $\delta'$ is fixed,
the constants below depend only on $\eta$. Let $r$ be the local
radius supplied by these lemmas. With
$T=\exp(\sqrt{\log x})$, we have
\[
    \lambda_T=\frac{c_0}{2\log(T+2)}
    \asymp \frac1{\sqrt{\log x}}.
\]
Thus $\rho=1/\log x<\lambda_T$ and $\rho<r/2$ for all sufficiently
large $x$.

By the definition of $I_H$ in~\eqref{eq:I-H},
\[
    I_H
    =
    \frac{1}{2\pi i}
    \int_{\mathcal H_{\mathrm{loc}}}
    F(s,y)\frac{x^s}{s}\,ds.
\]
Using the local factorization recorded in
Lemma~\ref{lem:G-local-analytic}, namely
\[
    F(s,y)=G(s,y)(s-1)^{-\alpha},
\]
and setting $s=1+z$, we obtain
\[
    I_H
    =
    \frac{x}{2\pi i}
    \int_{\mathcal H_z}
    \frac{G(1+z,y)}{1+z}
    z^{-\alpha}e^{z\log x}\,dz,
\]
where $\mathcal H_z$ is the image of $\mathcal H_{\mathrm{loc}}$
under $z=s-1$. Thus $\mathcal H_z$ is the truncated Hankel contour
whose arms run along the negative real axis from $-\lambda_T$ to
$-\rho$, with the circle $|z|=\rho$ oriented counterclockwise
through the right side of the origin.

For $z$ on $\mathcal H_z$, Lemma~\ref{lem:G-local-values} gives
\[
    G(1+z,y)=G(1,y)+O_\eta(|z|).
\]
Also, since $|z|\leq\lambda_T<1/2$ for large $x$,
\[
    \frac1{1+z}=1+O(|z|).
\]
Using the uniform boundedness of $G(1,y)$, we obtain
\begin{equation}\label{eq:hankel-local-error-factor}
    \frac{G(1+z,y)}{1+z}
    =G(1,y)+E(z,y),
    \qquad E(z,y)=O_\eta(|z|),
\end{equation}
uniformly for $|y|\leq\eta$ and $z\in\mathcal H_z$.

Accordingly, write
\[
    I_H=M_H+R_H,
\]
where
\[
    M_H
    :=
    \frac{xG(1,y)}{2\pi i}
    \int_{\mathcal H_z}z^{-\alpha}e^{z\log x}\,dz
\]
and
\[
    R_H
    :=
    \frac{x}{2\pi i}
    \int_{\mathcal H_z}E(z,y)z^{-\alpha}e^{z\log x}\,dz.
\]

We first evaluate $M_H$. Under the change of variables
$w=z\log x$, the contour $\mathcal H_z$ is mapped to a truncated
Hankel contour $\mathcal H_w$ in the $w$-plane, whose arms run
from $-\lambda_T\log x$ to $-1$ and whose circular part is
$|w|=1$. Then
\[
    M_H
    =
    xG(1,y)(\log x)^{\alpha-1}
    \frac1{2\pi i}
    \int_{\mathcal H_w}e^w w^{-\alpha}\,dw.
\]
By the standard Hankel representation of $1/\Gamma(\alpha)$~\textup{[]},
\[
    \frac1{2\pi i}
    \int_{\mathcal H_w}e^w w^{-\alpha}\,dw
    =
    \frac1{\Gamma(\alpha)}
    +O_\eta\left(e^{-\lambda_T\log x/2}\right).
\]
Therefore
\[
    M_H
    =
    xG(1,y)\frac{(\log x)^{\alpha-1}}{\Gamma(\alpha)}
    +O_\eta\left(
        x(\log x)^{\Real(\alpha)-1}
        e^{-\lambda_T\log x/2}
    \right).
\]
Since $\lambda_T\log x\asymp\sqrt{\log x}$, this tail error is
absorbed into $O_\eta(x(\log x)^{2\eta})$.

It remains to estimate $R_H$. From~\eqref{eq:hankel-local-error-factor},
\[
    |R_H|
    \ll_\eta
    x\int_{\mathcal H_z}
    |z|\,|z^{-\alpha}|\,e^{\Real(z)\log x}\,|dz|.
\]
On the Hankel contour, $|\operatorname{Arg}z|\leq\pi$ and
$|\Im\alpha|\leq2\eta$, so
\[
    |z^{-\alpha}|\ll_\eta |z|^{-\Real(\alpha)}.
\]
Thus
\[
    |R_H|
    \ll_\eta
    x\int_{\mathcal H_z}
    |z|^{1-\Real(\alpha)}e^{\Real(z)\log x}\,|dz|.
\]
Changing variables again by $w=z\log x$ gives
\[
    |R_H|
    \ll_\eta
    x(\log x)^{\Real(\alpha)-2}
    \int_{\mathcal H_w}
    |w|^{1-\Real(\alpha)}e^{\Real w}\,|dw|.
\]
The last integral is $O_\eta(1)$: this is immediate on $|w|=1$,
and on the arms it is bounded by
\[
    \int_1^{\lambda_T\log x} u^{1-\Real(\alpha)}e^{-u}\,du
    \leq
    \int_1^\infty u^{1-\Real(\alpha)}e^{-u}\,du
    \ll_\eta 1.
\]
Consequently,
\[
    R_H=O_\eta\left(x(\log x)^{\Real(\alpha)-2}\right).
\]
Since $\alpha=2+2y$, we have
\[
    \Real(\alpha)-2=2\Real(y)\leq2\eta,
\]
and hence
\[
    R_H=O_\eta(x(\log x)^{2\eta}).
\]

Combining the estimates for $M_H$ and $R_H$, and using
Lemma~\ref{lem:G-local-values} to replace $G(1,y)$ by $H(1,y)$,
we obtain
\[
    I_H
    =
    xH(1,y)\frac{(\log x)^{1+2y}}{\Gamma(2+2y)}
    +O_\eta(x(\log x)^{2\eta}),
\]
uniformly for $|y|\leq\eta$.
\end{proof}

